\section{回帰 : 関係を直線で表し, 予測に使う}

\subsection{問い : 広告費を増やしたら, 売上はどのくらい変わるか}

前章では, 2つの変数の関係の方向と強さを相関係数で測る方法を学んだ.
広告費と売上の相関係数は$r = 0.87$だった.
強い正の相関がある, つまり広告費が多い月は売上も多い傾向がある, ということまではわかった.

しかし, 知りたいのはその先だ.
広告費を10万円増やしたら, 売上はどのくらい変わるのか.
来月の広告費が決まったとき, 売上をどのくらいと見込めばよいのか.
相関係数はこの問いに答えてくれない.
$r = 0.87$は「強く連動している」とは教えてくれるが, 「いくつ変わるか」は教えてくれない.

相関係数が教えてくれないのは\textbf{傾き}だ.
散布図の点の並びに直線を当てはめれば, 「$x$が1増えたとき$y$が平均的にいくつ変わるか」を数値で示せる.
この直線を使えば, まだ見ていない$x$の値に対して$y$を予測することもできる.

本章を読み終えると, 散布図に「最も良い」直線を引き, その傾きと切片を解釈し, 直線の当てはまり具合を$R^2$で評価できるようになる.
さらに, 直線を引いたあとに残る\textbf{残差}を確認することで, 「この直線は信用してよいか」を自分で判断できるようになる.


\subsection{素朴な方法と限界 : 目で引いた線は当てになるか}

散布図に直線を引くこと自体は, 定規さえあればできる.
右上がりに並んだ点を眺めて, 「だいたいこのあたりかな」と線を引く.
日常の場面では, 実際にこれで済ませていることも多いだろう.

では, この方法はどこまで信用できるのか.

同じ散布図を3人に見せて, それぞれ「最も良い」と思う直線を引いてもらう.
3人の線はまず一致しない.
ある人は急な傾きの線を引き, 別の人はゆるやかな線を引く.
右端の点を重視する人もいれば, 中央の密集した点に引きずられる人もいる.

傾きが1.5倍も違う線が出てくることは珍しくない.
「広告費10万円増で売上5万円増」と読む人と, 「売上8万円増」と読む人が, 同じデータから出てくるのだ.
どちらの線が「正しい」のか, 目分量では決められない.

つまり, 2つの問題がある.
\textbf{基準がない}ことと, \textbf{再現性がない}ことだ.
「良い直線」とは何かの定義がなければ, 誰が引いても同じ線にはならない.
同じ人がもう一度引いても, 微妙に違う線になるかもしれない.

目分量に頼らず, 「良い直線」を一意に決める基準が必要だ.


\subsection{直線の当てはめ : 最小二乗法}

\subsubsection*{回帰直線とは}

2変数$x$と$y$の間に直線的な関係があると想定して,
\[
  y = a + bx
\]
という形でその関係を表すことを考える.
$b$は\textbf{傾き}で, $x$が1増えたときに$y$が平均的にいくつ変わるかを表す.
$a$は\textbf{切片}で, $x = 0$のときの$y$の値だ.
ただし, $x = 0$が現実的な意味を持たない場合も多い（広告費ゼロの月が存在しないなど）.
そのときは切片の数値自体にこだわる必要はない.

問題は, $a$と$b$をどう決めるかだ.

\subsubsection*{残差 : 直線からのずれ}

散布図のすべての点が1本の直線の上に乗ることは, ふつうない.
直線から外れた分, つまり「実際の値」と「直線が予測する値」のずれを\textbf{残差}と呼ぶ.

$i$番目のデータ点$(x_i, y_i)$に対して, 直線が予測する値は$a + bx_i$だ.
残差$e_i$は
\[
  e_i = y_i - (a + bx_i)
\]
で定義される.
残差が正なら実際の値が直線より上にあり, 負なら下にある.

すべての点を通る直線は引けない以上, 残差をゼロにすることはできない.
ならば, 残差をなるべく小さくする直線が「良い直線」だろう.

しかし, 残差をそのまま合計しても, 正の残差と負の残差が打ち消し合ってしまう.
合計がゼロに近いからといって, 直線がデータによく合っているとは限らない.

そこで, 各残差を\textbf{二乗}してから合計する.
二乗すれば正負の打ち消しが起きない.
この「残差の二乗和」を最小にする$a$と$b$を選ぶ方法が, \textbf{最小二乗法}だ.

\subsubsection*{最小二乗法 : 残差の二乗和を最小にする}

\begin{defi}{最小二乗法}{ols}
$n$個のデータ$(x_1, y_1), \ldots, (x_n, y_n)$に対し,
\[
  S(a, b) = \sum_{i=1}^{n} e_i^2 = \sum_{i=1}^{n} (y_i - a - bx_i)^2
\]
を最小にする$a$と$b$を選ぶ方法を\textbf{最小二乗法}（OLS: Ordinary Least Squares）と呼ぶ.
\end{defi}

$S(a, b)$は残差の二乗和だ.
これを最小にする$a$と$b$は, 計算によって一意に求まる.
結果は以下の通りである.

\[
  b = \frac{\displaystyle\sum_{i=1}^{n}(x_i - \bar{x})(y_i - \bar{y})}{\displaystyle\sum_{i=1}^{n}(x_i - \bar{x})^2}, \qquad a = \bar{y} - b\bar{x}
\]

傾き$b$の式を見てみよう.
分子は$x$と$y$の偏差の積の合計で, 前章で学んだ共分散の分子と同じ構造だ.
分母は$x$の偏差の二乗和で, $x$の分散の分子に対応する.
つまり傾き$b$は, $x$と$y$の共変動を$x$自身の変動で割ったものだ.
$x$が1単位ばらつくごとに$y$が平均的にどれだけ動くか, という直感に合っている.

切片$a$は$\bar{y} - b\bar{x}$で求まる.
これは, 回帰直線が必ず$(\bar{x}, \bar{y})$を通ることを意味している.
データの重心を通る直線を, 傾き$b$で引いている, と読める.

「なぜ二乗なのか, 絶対値ではだめなのか」と思うかもしれない.
絶対値を使う方法（最小絶対偏差法）も存在するし, 外れ値に対してはむしろ頑健だ.
ただし, 二乗のほうが数学的に扱いやすく, 解が公式で求まるという実用上の利点がある.
1809年にガウスが小惑星ケレスの軌道計算に最小二乗法を用いたとき, この「公式で解ける」性質が決定的に重要だった.

\begin{mybox}{ガウスとケレスの軌道}
1801年, イタリアの天文学者ピアッツィが小惑星ケレスを発見したが, 観測できたのはわずか41日間だった.
ケレスは太陽の向こう側に隠れてしまい, 再び姿を現したときにどこにいるかを予測する必要があった.
ガウスは限られた観測データから軌道を推定するために, 観測値と予測値のずれの二乗和を最小にする方法を用いた.
彼の予測は的中し, ケレスは予測された位置で再発見された.
最小二乗法は, このように「不完全なデータから最も確からしい答えを引き出す」という実務的な動機から生まれた道具だ.
なお, フランスのルジャンドルが1805年に最小二乗法を先に公表しており, 優先権を巡る論争がある.
ガウスは「自分は1795年から使っていた」と主張したが, 出版が後だったため決着はついていない.
\end{mybox}


\subsubsection*{相関係数との関係}

傾き$b$の式と相関係数$r$の式を見比べると, 次の関係があることがわかる.
\[
  b = r \cdot \frac{s_y}{s_x}
\]
$r$は関係の強さと方向を$-1$から$1$で表していた.
$b$はそれに$s_y / s_x$, つまり$y$と$x$の散らばりの比を掛けたものだ.
相関係数は「どれだけ一緒に動くか」を単位なしで測り, 傾きは「$x$が1増えたときに$y$がいくつ変わるか」を単位つきで測っている.
同じ関係を, 異なる角度から見ている.


\subsection{当てはまりの評価 : $R^2$と残差}

\subsubsection*{$R^2$ : 説明できた割合}

直線を引いた.
では, この直線はデータをどのくらいうまく説明しているのだろうか.

$y$の値は, データごとにばらつく.
そのばらつきの一部は, $x$との関係で説明できる.
残りは, 直線では説明しきれないばらつきだ.

\textbf{決定係数}$R^2$は, $y$の全体のばらつきのうち, 回帰直線で説明できた割合を表す.

\begin{defi}{決定係数}{r-squared}
$y$の全変動を$\mathrm{SS_{total}} = \sum(y_i - \bar{y})^2$, 回帰による変動を$\mathrm{SS_{reg}} = \sum(\hat{y}_i - \bar{y})^2$, 残差による変動を$\mathrm{SS_{res}} = \sum(y_i - \hat{y}_i)^2$とする.
決定係数は
\[
  R^2 = 1 - \frac{\mathrm{SS_{res}}}{\mathrm{SS_{total}}}
\]
で定義される. $\hat{y}_i = a + bx_i$は直線による予測値である.
\end{defi}

$R^2 = 0.87$なら, $y$のばらつきの87\%が$x$との直線関係で説明でき, 残り13\%は直線では説明できない, という意味だ.
$R^2 = 1$なら全点が直線上にあり, $R^2 = 0$なら$x$と$y$に線形関係がまったくない.

単回帰では, $R^2$は相関係数$r$の二乗に一致する.
$r = 0.87$なら$R^2 = 0.76$だ.
相関の「強さ」と回帰の「説明力」が, こうしてつながっている.

\subsubsection*{$R^2$だけでは安心できない}

$R^2$が高ければ, 直線がデータをよく捉えていると思いたくなる.
しかし, $R^2$は全体の当てはまりを1つの数値に圧縮した指標であり, ずれ方のパターンは見えない.

アンスコムのカルテットという有名な例がある.
4組のデータがあり, いずれも$R^2 \approx 0.67$, 傾きも切片もほぼ同じだ.
しかし散布図を描くと, 1組は素直な直線関係, 1組は曲線関係, 1組はほぼ完璧な直線に外れ値が1点, 1組は$x$が1点だけ離れていて見かけの相関を作っている.

% [図挿入: アンスコムのカルテットの4枚の散布図]

$R^2$の数値だけを見ていたら, 4組とも「まあまあの当てはまり」で片づけてしまう.
直線が妥当かどうかを判断するには, $R^2$に加えて\textbf{残差}を確認する必要がある.

\subsubsection*{残差プロット : 直線の妥当性を目で確かめる}

残差$e_i = y_i - \hat{y}_i$を, 予測値$\hat{y}_i$に対してプロットする.
この図が\textbf{残差プロット}だ.

直線がデータの関係をうまく捉えていれば, 残差はゼロのまわりにランダムに散らばる.
特定のパターンは見えないはずだ.

逆に, 残差に系統的なパターンが見えたら, 直線モデルでは捉えきれない構造がデータに残っている.
たとえば, 残差が予測値の大きさに応じてU字型に並んでいれば, 曲線的な関係を直線で無理に近似してしまった可能性がある.
残差の散らばりが予測値に応じて広がっていくなら, $x$の値が大きい領域で予測が不安定になっていることを示唆する.

残差プロットの読み方は2点に絞ればよい.
\begin{itemize}
  \item \textbf{パターンがあるか}. U字型, 扇形, 周期的なうねりなどが見えたら, 直線モデルの前提を疑う.
  \item \textbf{飛び離れた点があるか}. 他の残差から大きく外れた点は, 外れ値としてデータの確認が必要だ.
\end{itemize}

パターンがなく, ランダムに散らばっていれば, 「この直線はデータの関係をおおむね捉えている」と判断してよい.
厳密な検定もあるが, 本書の範囲では目視で十分だ.


\subsection{手法 : 回帰直線を求めて評価する}

\subsubsection*{前提条件}
\begin{itemize}
  \item 2変数の間に\textbf{直線的な関係}があること（散布図で確認する）.
  \item 各データ点が互いに\textbf{独立}であること.
  \item 極端な外れ値がないこと（あるいは, その影響を把握していること）.
\end{itemize}

\subsubsection*{手順}

\textbf{Step 1. 散布図を描く.}
まず$x$と$y$の散布図を描き, 直線的な関係がありそうかを目で確認する.
曲線的な関係しか見えないなら, 直線回帰をそのまま当てはめるのは不適切だ.

\textbf{Step 2. 傾き$b$と切片$a$を計算する.}
\[
  b = \frac{\sum(x_i - \bar{x})(y_i - \bar{y})}{\sum(x_i - \bar{x})^2}, \qquad a = \bar{y} - b\bar{x}
\]

\textbf{Step 3. $R^2$を計算する.}
\[
  R^2 = 1 - \frac{\sum(y_i - \hat{y}_i)^2}{\sum(y_i - \bar{y})^2}
\]

\textbf{Step 4. 残差プロットを描いて確認する.}
残差$e_i = y_i - \hat{y}_i$を予測値$\hat{y}_i$に対してプロットし, 系統的なパターンがないかを確認する.


\subsection{実践 : 広告費と売上で手を動かす}

ある店舗の10ヶ月分のデータを使う.

% TODO: 具体的な数値テーブルはランニング例確定後に挿入
% 以下は仮のデータで計算手順を示す

\subsubsection*{データと散布図}

% [表挿入: 月, 広告費(万円), 売上(万円) の10行]
% [図挿入: 広告費 vs 売上の散布図に回帰直線を重ねたもの]

散布図を描くと, 右上がりの傾向が見える.
直線を当てはめることに大きな無理はなさそうだ.

\subsubsection*{傾きと切片の計算}

Step 2の式に値を代入して計算する.

% TODO: 具体的な計算を仮データで実施

仮に$b = 2.3$, $a = 120$が得られたとしよう.
傾き$b = 2.3$は, 「広告費を1万円増やすと, 売上が平均的に2.3万円増える」と読める.
切片$a = 120$は, 「広告費がゼロのとき売上は120万円」という意味だが, 広告費ゼロの月が実際にあるかどうかはデータを見て判断する.
データの範囲外に直線を延長して解釈するのは危険だ.

\subsubsection*{$R^2$と残差の確認}

$R^2 = 0.76$だったとする.
売上のばらつきの76\%が広告費との直線関係で説明でき, 残り24\%は他の要因（季節, 天候, 競合の動きなど）によるものだ.

残差プロットを描く.

% [図挿入: 残差プロット]

残差がゼロのまわりにランダムに散らばっていれば, 直線モデルはおおむね妥当だと判断できる.
もしU字型のパターンが見えたら, 広告費と売上の関係は直線ではなく曲線的かもしれないし, ある広告費の水準を超えると効果が頭打ちになっている可能性もある.

\subsubsection*{回帰直線で「できること」と「できないこと」}

この直線を使えば, 「来月の広告費を50万円にしたとき, 売上はおよそ$120 + 2.3 \times 50 = 235$万円くらいだろう」と予測できる.
ただし, これはあくまでデータの範囲内での予測だ.
広告費が今までの範囲（たとえば20万円〜80万円）に収まっているなら, この予測にはそれなりの根拠がある.
広告費200万円のように, データの範囲をはるかに超えた値を代入して予測するのは, 直線関係がそこまで続いている保証がないため, 信頼性が低い.

もう一つ, 第7章で繰り返し確認したことを, ここでも改めて述べておく.
「広告費と売上に強い線形関係がある」ことと, 「広告費が売上の原因である」ことは別の話だ.
回帰直線は関係を記述し予測に使う道具であって, 因果を証明する道具ではない.


\subsection{結果と次の一手 : 直線を引いた先に残る問い}

前章では, 2変数の関係の方向と強さを相関係数で測った.
本章では, その関係を$y = a + bx$という直線で表し, 傾き$b$で「$x$が1増えたときの$y$の変化量」を数値にした.
最小二乗法により, 誰が計算しても同じ「残差の二乗和が最小になる直線」が得られる.
$R^2$で当てはまりを評価し, 残差プロットで直線の妥当性を目視で確認した.

これで, 「関係がある」と言うだけでなく, 「どのくらい変わるか」を定量的に語れるようになった.
予測にも使える.

ただし, ここまでの道具には前提がある.
データが独立に集められていること.
$x$と$y$の関係が直線的であること.
極端な外れ値がないこと.
これらの前提が崩れていたら, 直線を引いても数字は出るが, その数字は当てにならない.

前提が崩れる典型的な場面がある.
標本の集め方に偏りがある場合.
時間的に隣り合うデータが独立でない場合.
因果の方向を取り違えている場合.
こうした問題は, 直線を引く技術の外側にあり, データの設計と解釈の問題だ.

次章では, 2〜8章で積み上げてきた手法の前提を整理し, その前提が崩れるとどうなるかを体系的に確認する.
バイアスの類型, 前提の確認方法, そしてデータを扱うときの倫理.
手法の「外側」を知ることが, 手法を正しく使うための最後の一歩になる.
